\par\textbf{Do earlier returns help explain this week's return?} The Ljung--Box test checks whether a series is correlated with its own returns over the previous 13 weeks, roughly three months. It finds evidence of a relationship in 5 of 14 series: BCOMCOT, GBPUSD, SPXT, SX5T, TUKXG. This suggests that past returns may contain some information about subsequent returns. The test alone does not tell us how large that relationship is, whether it reflects continuation or reversal, or whether it is useful for forecasting.

\par\textbf{Does the size of market moves vary in persistent spells?} The ARCH test and the Ljung--Box test on squared, demeaned returns focus on the size of moves, regardless of whether prices rise or fall. They find evidence of dependence in 14 of 14 and 14 of 14 series, respectively, using 13 weekly lags. This is consistent with volatility clustering: turbulent weeks tend to occur in spells, as do calmer weeks. It supports allowing the model's estimate of risk to change over time. It does not, by itself, show that correlations between different markets change.

\par\textbf{What does this mean for our model?} We allow volatility to vary and will assess whether adding a simple dependence on last week's return improves the mean equation. These checks describe the observed data before model fitting. After estimation, we must repeat them on the part of returns the model leaves unexplained (the innovations) to see whether patterns remain. Dependence across weeks is not the same as a unit root and does not by itself call for differencing returns again.

\par{\footnotesize Reading the statistics: a rejection means the test found evidence against its no-dependence assumption at the chosen 5\% threshold. Testing many series creates more opportunities for chance flags; these counts have not been adjusted for that. Source: saved \texttt{diagnostics.csv}, percentage weekly log returns, 2003-08-15--2026-09-04; individual statistics and $p$-values remain in that file.}
